% 小行星卫星（综述）
% license CCBYSA3
% type Wiki

本文根据 CC-BY-SA 协议转载翻译自维基百科\href{https://en.wikipedia.org/wiki/Minor-planet_moon}{相关文章}。

\begin{figure}[ht]
\centering
\includegraphics[width=6cm]{./figures/fe309a0bf64922bd.png}
\caption{} \label{fig_Minor_2}
\end{figure}
一个小行星卫星是指作为天然卫星环绕一颗小行星运行的天文对象。截至 2022 年 1 月，已有 457 颗已知或被怀疑拥有卫星的小行星。之所以说对小行星卫星（以及一般意义上的双天体系统）的发现十分重要，是因为通过确定它们的轨道，可以估计主星的质量与密度，从而获得通常难以通过其他方式获取的物理性质信息。

在这些卫星中，有若干在相对尺寸上非常巨大：90 号 Antiope、Mors–Somnus 与 Sila–Nunam 的尺寸约为其主星的 95\%；Patroclus–Menoetius、Altjira 与 Lempo–Hiisi 的尺寸约为 90\%，其中 Lempo–Paha 约为 50\%。已知绝对尺寸最大的小行星卫星是冥王星的卫星 Charon，其直径约为冥王星的一半。

此外，在若干遥远的天体周围，人们还发现了多个环系统（参见：Chariklo 与 Chiron 的环）。
\subsubsection{术语}
除了使用“卫星”和“卫星（月）”这两个术语之外，“双星”（双小行星）一词有时也用于指拥有一颗卫星的小行星，而“三级系统”则用于指拥有两颗卫星的小行星。如果其中一个天体显著更大，则称其为主星（primary），而其伴随天体称为伴星（secondary）。当小行星及其卫星的尺寸大致相当时，有时使用“双小行星”一词，而“binary”这一术语通常与两者的相对大小无关。当双小行星系统的两个天体大小相近时，小行星中心（MPC）将它们称为“binary companions”，而不是将较小的天体称为卫星。一个真正双星系统的良好示例是 90 号 Antiope 系统，其性质于 2000 年 8 月确认。非常小的卫星通常被称为“moonlets”。
\subsection{发现历程}
在哈勃太空望远镜时代和深空探测器抵达太阳系外部区域之前，寻找小行星是否具有卫星的努力主要受限于地球上的光学观测。例如，在 1978 年，通过恒星掩星观测，人们曾声称检测到 532 号 Herculina 存在一颗卫星。然而，后来哈勃太空望远镜的更高分辨率成像并未发现任何卫星，目前的共识是 Herculina 并不具有显著的卫星。在随后的几年中，也出现过其他类似的小行星伴星（通常称作卫星）观测报告。当时，天文学家 Thomas Hamilton 在《Sky \& Telescope》杂志发表的一封信中指出，地球上存在一些看似同时形成的撞击坑（例如魁北克的 Clearwater 湖），这些撞击坑可能是由成对、相互引力束缚的天体造成的。

同样在 1978 年，冥王星最大的卫星 Charon 被发现；但在当时，冥王星仍被视为一颗主要行星。

1993 年，当伽利略号探测器发现小天体 Dactyl 环绕小行星 243 号 Ida 运行时，第一个被确认的小行星卫星得以成立。第二例卫星于 1998 年在小行星 45 号 Eugenia 周围被发现。2001 年，小行星 617 Patroclus 及其等尺寸伴星 Menoetius 成为木星特洛伊群中首对已知的双小行星系统。冥王星–Charon 系统之后，第一个被光学分辨的海王星外双星——1998 WW31——于 2002 年获得确认。

2021 年，人们发现 130 号 Elektra 拥有三颗卫星，使其成为目前唯一已知的四重小行星系统。
